\section{Measure the complete conversion boundary}
Count simulated instructions from function entry to its return boundary, including
preflight, frame handling, validation and any kernel calls. Exclude caller setup,
assembly and session creation. The following measurements use ordinary code/data addresses and the same binary pattern
\((73i+129)\bmod256\) for both paths. Decoder input is the canonical encoding of that
binary pattern, so its input count differs from the raw count in the first column.
\tagpdfsetup{table/header-rows={1}}
\begin{center}
  \begin{tabular}{rrrrr}
    Raw bytes & Encode scalar & Encode MSA & Decode scalar & Decode MSA \\
    \hline
    12        & 244           & 178        & 1078          & 675        \\
    48        & 820           & 562        & 4196          & 2510       \\
    192       & 3124          & 2098       & 16496         & 9764       \\
    4096      & 65598         & 43778      & 349870        & 206555     \\
  \end{tabular}
\end{center}

The direct encoder kernel executes 110 instructions; the direct decoder kernel executes
134. Those kernel counts exclude whole-buffer validation. Comparing only a kernel with
a complete scalar decoder would therefore use different boundaries. The alternative
scalar implementation also changes character classification from branches to a table,
which changes the cost of both decoder passes independently of vector packing.
For 5464 text bytes, its scalar path takes 357886 instructions, compared with 349870
for the branch-based path on this pattern. A table lookup requires address formation
and a byte load; count actual instructions rather than assuming a lookup is faster.

These measurements include the cumulative starter’s unfinished-call guards and shared
restore branches, even after every region has been completed. The guards preserve a
useful development boundary but still execute in the completed program.

These counts describe current simulated work. They do not measure host latency or
hardware cycles, and no exact count is a grading rule. Data can affect the branch-based
classifier's path. Record the actual input pattern and whether the last quartet was
padded when making another comparison.

As an optimization exercise, investigate replacing two vector range comparisons with
a subtraction followed by one unsigned width test. Prove the byte interval and wraparound
behavior, and test all 256 byte values before using the new masks. Also investigate
an aligned three-word copy fast path. Check the low two address bits before word loads,
and retain byte copies for unaligned input. Both paths must consume exactly twelve
source bytes.
